\documentclass[10pt,letterpaper]{article}
\usepackage{cornell-notes}
\usepackage{mathtools}

\title{Chapter 2: Solution of Linear Algebraic Equations}
\author{}
\date{}

\setDocTitle{Chapter 2: Solution of Linear Algebraic Equations}
\setDocSubtitle{Numerical Methods}
\setDocVersion{v1.0}
\setDocStatus{Study Companion}
\setDocOwner{Numerical Methods Cornell Notes}
\setDocAuthor{}
\setDocDate{\today}
\setDocSummary{Cornell-format study and review notes for Chapter 2: Solution of Linear Algebraic Equations.}

\setCornellCollection{Numerical Methods Cornell Notes}
\setCornellUnitType{Chapter}
\setCornellUnitNumber{2}
\setCornellUnitTitle{Solution of Linear Algebraic Equations}
\setCornellCompanionLabel{Study and review companion}

\hypersetup{pdftitle={Chapter 2: Solution of Linear Algebraic Equations},pdfsubject={Cornell notes},pdfauthor={}}

\begin{document}
\maketitle

\section*{Chapter Overview}
This chapter organizes the mathematical and computational ideas behind solution of linear algebraic equations. The notes preserve the numbered section sequence while emphasizing method selection, explicit assumptions, numerical reliability, cost, and verification.

\section*{Learning Objectives}
Explain the governing problem class; compare Gauss-Jordan Elimination, Gaussian Elimination with Backsubstitution, LU Decomposition and Its Applications, and Tridiagonal and Band-Diagonal Systems of Equations; design defensible stopping and verification checks; distinguish problem conditioning from algorithmic stability.

\section*{Prerequisites}
Matrix multiplication, triangular systems, norms, and floating-point error.

\section*{Operational Workflow}
Inspect matrix shape, symmetry, definiteness, sparsity, scaling, and reuse; select a factorization or iterative method; solve; then verify with residual and sensitivity information.

\subsection*{Formula and notation anchor}
\[
A\widehat{x}=b-r,\qquad r=b-A\widehat{x},\qquad \frac{\|\widehat{x}-x\|}{\|x\|}\lesssim \kappa(A)\frac{\|r\|}{\|b\|}
\]
A small residual implies a small forward error only when the matrix is not badly conditioned.

\begin{warnbox}{Numerical warning}
Small pivots, ill-conditioning, fill-in, rank deficiency, misleadingly small residuals, and explicit inversion when a direct solve is sufficient.
\end{warnbox}

\section{2.0 Introduction}
\begin{cornell}
\noterow{What is the central idea?}{Linear systems connect algebraic models to computation; method choice depends on structure, conditioning, sparsity, and whether one or many right-hand sides are required.}
\noterow{How should it be used?}{For Introduction, begin from explicit inputs, outputs, preconditions, and representation choices.  Inspect matrix shape, symmetry, definiteness, sparsity, scaling, and reuse; select a factorization or iterative method; solve; then verify with residual and sensitivity information.}
\noterow{Cost and reliability}{Analyze arithmetic work, storage, data movement, and convergence separately.  Relevant failure pressures include small pivots, ill-conditioning, fill-in, rank deficiency, misleadingly small residuals, and explicit inversion when a direct solve is sufficient.  Use scaled tolerances and report both success criteria and diagnostic evidence.}
\noterow{Implementation checklist}{\begin{itemize}
\item Validate dimensions, domains, ordering, and structural assumptions.
\item Guard small divisors, invalid states, overflow, underflow, and nonfinite values.
\item Make mutation, workspace, indexing, and termination behavior explicit.
\item Test a simple exact case, a difficult valid case, a boundary case, and an expected failure.
\end{itemize}}
\end{cornell}
\begin{summarybox}{Section summary}
Linear systems connect algebraic models to computation; method choice depends on structure, conditioning, sparsity, and whether one or many right-hand sides are required.  Select this technique only when its assumptions match the data, and accept a result only after an independent residual, error estimate, invariant, or refinement check.
\end{summarybox}

\section{2.1 Gauss-Jordan Elimination}
\begin{cornell}
\noterow{What is the central idea?}{Gauss-Jordan elimination applies row operations until the coefficient matrix becomes the identity, producing a solution and, when requested, an inverse.}
\noterow{How should it be used?}{For Gauss-Jordan Elimination, begin from explicit inputs, outputs, preconditions, and representation choices.  Inspect matrix shape, symmetry, definiteness, sparsity, scaling, and reuse; select a factorization or iterative method; solve; then verify with residual and sensitivity information.}
\noterow{Quantitative anchor}{\[
[A\mid b]\longrightarrow[I\mid x]
\]
State the dimensions, scaling, and validity assumptions before using this relation.  Check the resulting residual or invariant rather than trusting an iteration count alone.}
\noterow{Cost and reliability}{Analyze arithmetic work, storage, data movement, and convergence separately.  Relevant failure pressures include small pivots, ill-conditioning, fill-in, rank deficiency, misleadingly small residuals, and explicit inversion when a direct solve is sufficient.  Use scaled tolerances and report both success criteria and diagnostic evidence.}
\noterow{Implementation checklist}{\begin{itemize}
\item Validate dimensions, domains, ordering, and structural assumptions.
\item Guard small divisors, invalid states, overflow, underflow, and nonfinite values.
\item Make mutation, workspace, indexing, and termination behavior explicit.
\item Test a simple exact case, a difficult valid case, a boundary case, and an expected failure.
\end{itemize}}
\end{cornell}
\begin{summarybox}{Section summary}
Gauss-Jordan elimination applies row operations until the coefficient matrix becomes the identity, producing a solution and, when requested, an inverse.  Select this technique only when its assumptions match the data, and accept a result only after an independent residual, error estimate, invariant, or refinement check.
\end{summarybox}

\section{2.2 Gaussian Elimination with Backsubstitution}
\begin{cornell}
\noterow{What is the central idea?}{Forward elimination creates an upper-triangular system and back substitution recovers unknowns; pivoting protects the computation from small divisors.}
\noterow{How should it be used?}{For Gaussian Elimination with Backsubstitution, begin from explicit inputs, outputs, preconditions, and representation choices.  Inspect matrix shape, symmetry, definiteness, sparsity, scaling, and reuse; select a factorization or iterative method; solve; then verify with residual and sensitivity information.}
\noterow{Quantitative lens}{Identify the governing equation, recurrence, probability law, or geometric predicate for Gaussian Elimination with Backsubstitution.  Define every variable and scale; then derive a residual, error estimate, or invariant that can be checked independently.}
\noterow{Cost and reliability}{Analyze arithmetic work, storage, data movement, and convergence separately.  Relevant failure pressures include small pivots, ill-conditioning, fill-in, rank deficiency, misleadingly small residuals, and explicit inversion when a direct solve is sufficient.  Use scaled tolerances and report both success criteria and diagnostic evidence.}
\noterow{Implementation checklist}{\begin{itemize}
\item Validate dimensions, domains, ordering, and structural assumptions.
\item Guard small divisors, invalid states, overflow, underflow, and nonfinite values.
\item Make mutation, workspace, indexing, and termination behavior explicit.
\item Test a simple exact case, a difficult valid case, a boundary case, and an expected failure.
\end{itemize}}
\end{cornell}
\begin{summarybox}{Section summary}
Forward elimination creates an upper-triangular system and back substitution recovers unknowns; pivoting protects the computation from small divisors.  Select this technique only when its assumptions match the data, and accept a result only after an independent residual, error estimate, invariant, or refinement check.
\end{summarybox}

\section{2.3 LU Decomposition and Its Applications}
\begin{cornell}
\noterow{What is the central idea?}{LU factorization separates elimination from substitution, making repeated solves, determinant evaluation, and inverse-related operations more efficient.}
\noterow{How should it be used?}{For LU Decomposition and Its Applications, begin from explicit inputs, outputs, preconditions, and representation choices.  Inspect matrix shape, symmetry, definiteness, sparsity, scaling, and reuse; select a factorization or iterative method; solve; then verify with residual and sensitivity information.}
\noterow{Quantitative anchor}{\[
PA=LU,\qquad Ly=Pb,\quad Ux=y
\]
State the dimensions, scaling, and validity assumptions before using this relation.  Check the resulting residual or invariant rather than trusting an iteration count alone.}
\noterow{Cost and reliability}{Analyze arithmetic work, storage, data movement, and convergence separately.  Relevant failure pressures include small pivots, ill-conditioning, fill-in, rank deficiency, misleadingly small residuals, and explicit inversion when a direct solve is sufficient.  Use scaled tolerances and report both success criteria and diagnostic evidence.}
\noterow{Implementation checklist}{\begin{itemize}
\item Validate dimensions, domains, ordering, and structural assumptions.
\item Guard small divisors, invalid states, overflow, underflow, and nonfinite values.
\item Make mutation, workspace, indexing, and termination behavior explicit.
\item Test a simple exact case, a difficult valid case, a boundary case, and an expected failure.
\end{itemize}}
\end{cornell}
\begin{summarybox}{Section summary}
LU factorization separates elimination from substitution, making repeated solves, determinant evaluation, and inverse-related operations more efficient.  Select this technique only when its assumptions match the data, and accept a result only after an independent residual, error estimate, invariant, or refinement check.
\end{summarybox}

\section{2.4 Tridiagonal and Band-Diagonal Systems of Equations}
\begin{cornell}
\noterow{What is the central idea?}{Banded structure permits elimination that stores and updates only nearby diagonals, reducing dense cubic work to a cost proportional to bandwidth.}
\noterow{How should it be used?}{For Tridiagonal and Band-Diagonal Systems of Equations, begin from explicit inputs, outputs, preconditions, and representation choices.  Inspect matrix shape, symmetry, definiteness, sparsity, scaling, and reuse; select a factorization or iterative method; solve; then verify with residual and sensitivity information.}
\noterow{Quantitative lens}{Identify the governing equation, recurrence, probability law, or geometric predicate for Tridiagonal and Band-Diagonal Systems of Equations.  Define every variable and scale; then derive a residual, error estimate, or invariant that can be checked independently.}
\noterow{Cost and reliability}{Analyze arithmetic work, storage, data movement, and convergence separately.  Relevant failure pressures include small pivots, ill-conditioning, fill-in, rank deficiency, misleadingly small residuals, and explicit inversion when a direct solve is sufficient.  Use scaled tolerances and report both success criteria and diagnostic evidence.}
\noterow{Implementation checklist}{\begin{itemize}
\item Validate dimensions, domains, ordering, and structural assumptions.
\item Guard small divisors, invalid states, overflow, underflow, and nonfinite values.
\item Make mutation, workspace, indexing, and termination behavior explicit.
\item Test a simple exact case, a difficult valid case, a boundary case, and an expected failure.
\end{itemize}}
\end{cornell}
\begin{summarybox}{Section summary}
Banded structure permits elimination that stores and updates only nearby diagonals, reducing dense cubic work to a cost proportional to bandwidth.  Select this technique only when its assumptions match the data, and accept a result only after an independent residual, error estimate, invariant, or refinement check.
\end{summarybox}

\section{2.5 Iterative Improvement of a Solution to Linear Equations}
\begin{cornell}
\noterow{What is the central idea?}{Iterative refinement solves a correction equation from the residual, often recovering digits when residuals are accumulated more accurately than the factorization.}
\noterow{How should it be used?}{For Iterative Improvement of a Solution to Linear Equations, begin from explicit inputs, outputs, preconditions, and representation choices.  Inspect matrix shape, symmetry, definiteness, sparsity, scaling, and reuse; select a factorization or iterative method; solve; then verify with residual and sensitivity information.}
\noterow{Quantitative lens}{Identify the governing equation, recurrence, probability law, or geometric predicate for Iterative Improvement of a Solution to Linear Equations.  Define every variable and scale; then derive a residual, error estimate, or invariant that can be checked independently.}
\noterow{Cost and reliability}{Analyze arithmetic work, storage, data movement, and convergence separately.  Relevant failure pressures include small pivots, ill-conditioning, fill-in, rank deficiency, misleadingly small residuals, and explicit inversion when a direct solve is sufficient.  Use scaled tolerances and report both success criteria and diagnostic evidence.}
\noterow{Implementation checklist}{\begin{itemize}
\item Validate dimensions, domains, ordering, and structural assumptions.
\item Guard small divisors, invalid states, overflow, underflow, and nonfinite values.
\item Make mutation, workspace, indexing, and termination behavior explicit.
\item Test a simple exact case, a difficult valid case, a boundary case, and an expected failure.
\end{itemize}}
\end{cornell}
\begin{summarybox}{Section summary}
Iterative refinement solves a correction equation from the residual, often recovering digits when residuals are accumulated more accurately than the factorization.  Select this technique only when its assumptions match the data, and accept a result only after an independent residual, error estimate, invariant, or refinement check.
\end{summarybox}

\section{2.6 Singular Value Decomposition}
\begin{cornell}
\noterow{What is the central idea?}{The SVD resolves a matrix into orthogonal directions and nonnegative singular values, exposing rank, null spaces, least-squares structure, and sensitivity.}
\noterow{How should it be used?}{For Singular Value Decomposition, begin from explicit inputs, outputs, preconditions, and representation choices.  Inspect matrix shape, symmetry, definiteness, sparsity, scaling, and reuse; select a factorization or iterative method; solve; then verify with residual and sensitivity information.}
\noterow{Quantitative anchor}{\[
A=U\Sigma V^T,\qquad A^+=V\Sigma^+U^T
\]
State the dimensions, scaling, and validity assumptions before using this relation.  Check the resulting residual or invariant rather than trusting an iteration count alone.}
\noterow{Cost and reliability}{Analyze arithmetic work, storage, data movement, and convergence separately.  Relevant failure pressures include small pivots, ill-conditioning, fill-in, rank deficiency, misleadingly small residuals, and explicit inversion when a direct solve is sufficient.  Use scaled tolerances and report both success criteria and diagnostic evidence.}
\noterow{Implementation checklist}{\begin{itemize}
\item Validate dimensions, domains, ordering, and structural assumptions.
\item Guard small divisors, invalid states, overflow, underflow, and nonfinite values.
\item Make mutation, workspace, indexing, and termination behavior explicit.
\item Test a simple exact case, a difficult valid case, a boundary case, and an expected failure.
\end{itemize}}
\end{cornell}
\begin{summarybox}{Section summary}
The SVD resolves a matrix into orthogonal directions and nonnegative singular values, exposing rank, null spaces, least-squares structure, and sensitivity.  Select this technique only when its assumptions match the data, and accept a result only after an independent residual, error estimate, invariant, or refinement check.
\end{summarybox}

\section{2.7 Sparse Linear Systems}
\begin{cornell}
\noterow{What is the central idea?}{Sparse methods exploit a small set of nonzeros through suitable storage, matrix-vector products, preconditioning, and convergence monitoring rather than dense factorization.}
\noterow{How should it be used?}{For Sparse Linear Systems, begin from explicit inputs, outputs, preconditions, and representation choices.  Inspect matrix shape, symmetry, definiteness, sparsity, scaling, and reuse; select a factorization or iterative method; solve; then verify with residual and sensitivity information.}
\noterow{Quantitative lens}{Identify the governing equation, recurrence, probability law, or geometric predicate for Sparse Linear Systems.  Define every variable and scale; then derive a residual, error estimate, or invariant that can be checked independently.}
\noterow{Cost and reliability}{Analyze arithmetic work, storage, data movement, and convergence separately.  Relevant failure pressures include small pivots, ill-conditioning, fill-in, rank deficiency, misleadingly small residuals, and explicit inversion when a direct solve is sufficient.  Use scaled tolerances and report both success criteria and diagnostic evidence.}
\noterow{Implementation checklist}{\begin{itemize}
\item Validate dimensions, domains, ordering, and structural assumptions.
\item Guard small divisors, invalid states, overflow, underflow, and nonfinite values.
\item Make mutation, workspace, indexing, and termination behavior explicit.
\item Test a simple exact case, a difficult valid case, a boundary case, and an expected failure.
\end{itemize}}
\end{cornell}
\begin{summarybox}{Section summary}
Sparse methods exploit a small set of nonzeros through suitable storage, matrix-vector products, preconditioning, and convergence monitoring rather than dense factorization.  Select this technique only when its assumptions match the data, and accept a result only after an independent residual, error estimate, invariant, or refinement check.
\end{summarybox}

\section{2.8 Vandermonde Matrices and Toeplitz Matrices}
\begin{cornell}
\noterow{What is the central idea?}{Structured matrices admit specialized recurrences, but their apparent algebraic simplicity does not remove conditioning and cancellation risks.}
\noterow{How should it be used?}{For Vandermonde Matrices and Toeplitz Matrices, begin from explicit inputs, outputs, preconditions, and representation choices.  Inspect matrix shape, symmetry, definiteness, sparsity, scaling, and reuse; select a factorization or iterative method; solve; then verify with residual and sensitivity information.}
\noterow{Quantitative lens}{Identify the governing equation, recurrence, probability law, or geometric predicate for Vandermonde Matrices and Toeplitz Matrices.  Define every variable and scale; then derive a residual, error estimate, or invariant that can be checked independently.}
\noterow{Cost and reliability}{Analyze arithmetic work, storage, data movement, and convergence separately.  Relevant failure pressures include small pivots, ill-conditioning, fill-in, rank deficiency, misleadingly small residuals, and explicit inversion when a direct solve is sufficient.  Use scaled tolerances and report both success criteria and diagnostic evidence.}
\noterow{Implementation checklist}{\begin{itemize}
\item Validate dimensions, domains, ordering, and structural assumptions.
\item Guard small divisors, invalid states, overflow, underflow, and nonfinite values.
\item Make mutation, workspace, indexing, and termination behavior explicit.
\item Test a simple exact case, a difficult valid case, a boundary case, and an expected failure.
\end{itemize}}
\end{cornell}
\begin{summarybox}{Section summary}
Structured matrices admit specialized recurrences, but their apparent algebraic simplicity does not remove conditioning and cancellation risks.  Select this technique only when its assumptions match the data, and accept a result only after an independent residual, error estimate, invariant, or refinement check.
\end{summarybox}

\section{2.9 Cholesky Decomposition}
\begin{cornell}
\noterow{What is the central idea?}{A symmetric positive-definite matrix factors as a triangular factor times its transpose, providing a stable solve with roughly half the work of general LU.}
\noterow{How should it be used?}{For Cholesky Decomposition, begin from explicit inputs, outputs, preconditions, and representation choices.  Inspect matrix shape, symmetry, definiteness, sparsity, scaling, and reuse; select a factorization or iterative method; solve; then verify with residual and sensitivity information.}
\noterow{Quantitative anchor}{\[
A=LL^T\quad\text{for symmetric positive-definite }A
\]
State the dimensions, scaling, and validity assumptions before using this relation.  Check the resulting residual or invariant rather than trusting an iteration count alone.}
\noterow{Cost and reliability}{Analyze arithmetic work, storage, data movement, and convergence separately.  Relevant failure pressures include small pivots, ill-conditioning, fill-in, rank deficiency, misleadingly small residuals, and explicit inversion when a direct solve is sufficient.  Use scaled tolerances and report both success criteria and diagnostic evidence.}
\noterow{Implementation checklist}{\begin{itemize}
\item Validate dimensions, domains, ordering, and structural assumptions.
\item Guard small divisors, invalid states, overflow, underflow, and nonfinite values.
\item Make mutation, workspace, indexing, and termination behavior explicit.
\item Test a simple exact case, a difficult valid case, a boundary case, and an expected failure.
\end{itemize}}
\end{cornell}
\begin{summarybox}{Section summary}
A symmetric positive-definite matrix factors as a triangular factor times its transpose, providing a stable solve with roughly half the work of general LU.  Select this technique only when its assumptions match the data, and accept a result only after an independent residual, error estimate, invariant, or refinement check.
\end{summarybox}

\section{2.10 QR Decomposition}
\begin{cornell}
\noterow{What is the central idea?}{QR factorization separates a matrix into an orthogonal factor and an upper-triangular factor, supporting stable least-squares solves and eigenvalue algorithms.}
\noterow{How should it be used?}{For QR Decomposition, begin from explicit inputs, outputs, preconditions, and representation choices.  Inspect matrix shape, symmetry, definiteness, sparsity, scaling, and reuse; select a factorization or iterative method; solve; then verify with residual and sensitivity information.}
\noterow{Quantitative anchor}{\[
A=QR,\qquad Q^TQ=I
\]
State the dimensions, scaling, and validity assumptions before using this relation.  Check the resulting residual or invariant rather than trusting an iteration count alone.}
\noterow{Cost and reliability}{Analyze arithmetic work, storage, data movement, and convergence separately.  Relevant failure pressures include small pivots, ill-conditioning, fill-in, rank deficiency, misleadingly small residuals, and explicit inversion when a direct solve is sufficient.  Use scaled tolerances and report both success criteria and diagnostic evidence.}
\noterow{Implementation checklist}{\begin{itemize}
\item Validate dimensions, domains, ordering, and structural assumptions.
\item Guard small divisors, invalid states, overflow, underflow, and nonfinite values.
\item Make mutation, workspace, indexing, and termination behavior explicit.
\item Test a simple exact case, a difficult valid case, a boundary case, and an expected failure.
\end{itemize}}
\end{cornell}
\begin{summarybox}{Section summary}
QR factorization separates a matrix into an orthogonal factor and an upper-triangular factor, supporting stable least-squares solves and eigenvalue algorithms.  Select this technique only when its assumptions match the data, and accept a result only after an independent residual, error estimate, invariant, or refinement check.
\end{summarybox}

\section{2.11 Is Matrix Inversion an N\textasciicircum{}3 Process?}
\begin{cornell}
\noterow{What is the central idea?}{The cost of inversion depends on matrix structure and multiplication algorithms, but explicit inversion is usually inferior to solving the required systems directly.}
\noterow{How should it be used?}{For Is Matrix Inversion an N\textasciicircum{}3 Process?, begin from explicit inputs, outputs, preconditions, and representation choices.  Inspect matrix shape, symmetry, definiteness, sparsity, scaling, and reuse; select a factorization or iterative method; solve; then verify with residual and sensitivity information.}
\noterow{Quantitative lens}{Identify the governing equation, recurrence, probability law, or geometric predicate for Is Matrix Inversion an N\textasciicircum{}3 Process?.  Define every variable and scale; then derive a residual, error estimate, or invariant that can be checked independently.}
\noterow{Cost and reliability}{Analyze arithmetic work, storage, data movement, and convergence separately.  Relevant failure pressures include small pivots, ill-conditioning, fill-in, rank deficiency, misleadingly small residuals, and explicit inversion when a direct solve is sufficient.  Use scaled tolerances and report both success criteria and diagnostic evidence.}
\noterow{Implementation checklist}{\begin{itemize}
\item Validate dimensions, domains, ordering, and structural assumptions.
\item Guard small divisors, invalid states, overflow, underflow, and nonfinite values.
\item Make mutation, workspace, indexing, and termination behavior explicit.
\item Test a simple exact case, a difficult valid case, a boundary case, and an expected failure.
\end{itemize}}
\end{cornell}
\begin{summarybox}{Section summary}
The cost of inversion depends on matrix structure and multiplication algorithms, but explicit inversion is usually inferior to solving the required systems directly.  Select this technique only when its assumptions match the data, and accept a result only after an independent residual, error estimate, invariant, or refinement check.
\end{summarybox}

\section*{Method comparison}
\begin{center}
\small
\begin{tabularx}{\linewidth}{>{\bfseries\raggedright\arraybackslash}p{0.22\linewidth}XX}
\toprule
Method or lens & Best use & Main caution \\
\midrule
LU with pivoting & General dense square systems & $O(n^3)$ factorization; pivot growth \\
Cholesky & Symmetric positive-definite systems & Fast and stable; fails if assumptions are false \\
QR or SVD & Least squares or rank uncertainty & More work; stronger diagnostics \\
\bottomrule
\end{tabularx}
\end{center}

\section*{Worked example}
\begin{exambox}{Independent worked example}
Solve $\begin{bmatrix}4&1\\1&3\end{bmatrix}x=\begin{bmatrix}1\\2\end{bmatrix}$.  Elimination gives $x=(1/11,7/11)^T$.  Substitution verifies $Ax=b$ exactly in rational arithmetic.  In floating point, compute $r=b-A\widehat{x}$ and compare $\|r\|$ with $\|A\|\|\widehat{x}\|+\|b\|$ before accepting the result.
\end{exambox}

\section*{Chapter synthesis}
\begin{summarybox}{Chapter synthesis}
The chapter's methods should be viewed as a decision system rather than a list of routines.  Begin with the mathematical problem and its conditioning, exploit structure, select an algorithm with compatible stability and cost, and verify the computed answer with evidence that is independent of the internal iteration.  The recurring implementation discipline is: inspect matrix shape, symmetry, definiteness, sparsity, scaling, and reuse; select a factorization or iterative method; solve; then verify with residual and sensitivity information.
\end{summarybox}

\subsection*{Common mistakes and numerical hazards}
\begin{warnbox}{Common mistakes and numerical hazards}
Small pivots, ill-conditioning, fill-in, rank deficiency, misleadingly small residuals, and explicit inversion when a direct solve is sufficient.  A second recurring mistake is reporting only a computed value: always retain tolerances, iteration counts, residuals, refinement behavior, and relevant structural checks.
\end{warnbox}

\subsection*{Retrieval practice}
\textbf{Questions}
\begin{enumerate}
\item What problem does Gauss-Jordan Elimination address, and which assumption is easiest to violate?
\item What problem does Gaussian Elimination with Backsubstitution address, and which assumption is easiest to violate?
\item What problem does LU Decomposition and Its Applications address, and which assumption is easiest to violate?
\item What problem does QR Decomposition address, and which assumption is easiest to violate?
\item What problem does Is Matrix Inversion an N\textasciicircum{}3 Process? address, and which assumption is easiest to violate?
\item How do conditioning, stability, and the final residual play different roles in this chapter?
\end{enumerate}

\textbf{Brief answers}
\begin{enumerate}
\item Gauss-Jordan elimination applies row operations until the coefficient matrix becomes the identity, producing a solution and, when requested, an inverse. Recheck the section's domain, structure, scaling, and failure diagnostics.
\item Forward elimination creates an upper-triangular system and back substitution recovers unknowns; pivoting protects the computation from small divisors. Recheck the section's domain, structure, scaling, and failure diagnostics.
\item LU factorization separates elimination from substitution, making repeated solves, determinant evaluation, and inverse-related operations more efficient. Recheck the section's domain, structure, scaling, and failure diagnostics.
\item QR factorization separates a matrix into an orthogonal factor and an upper-triangular factor, supporting stable least-squares solves and eigenvalue algorithms. Recheck the section's domain, structure, scaling, and failure diagnostics.
\item The cost of inversion depends on matrix structure and multiplication algorithms, but explicit inversion is usually inferior to solving the required systems directly. Recheck the section's domain, structure, scaling, and failure diagnostics.
\item Conditioning describes sensitivity of the mathematical problem, stability describes error propagation by the algorithm, and the residual measures satisfaction of the computed equations; none is a universal substitute for the others.
\end{enumerate}

\subsection*{Mastery checklist}
\begin{itemize}
\item[$\square$] I can explain and select Gauss-Jordan Elimination without relying on a memorized code listing.
\item[$\square$] I can explain and select Gaussian Elimination with Backsubstitution without relying on a memorized code listing.
\item[$\square$] I can explain and select LU Decomposition and Its Applications without relying on a memorized code listing.
\item[$\square$] I can explain and select Tridiagonal and Band-Diagonal Systems of Equations without relying on a memorized code listing.
\item[$\square$] I can explain and select Iterative Improvement of a Solution to Linear Equations without relying on a memorized code listing.
\item[$\square$] I can state a scaled stopping rule and an independent verification check.
\item[$\square$] I can identify at least one conditioning risk and one algorithmic stability risk.
\item[$\square$] I can estimate time, storage, and data-movement costs at the appropriate level.
\end{itemize}

\end{document}

